\section{战败之后：改革为何成为地方政治}

1894年甲午战争爆发，1895年《马关条约》把赔款、通商与台湾割让推到清廷必须处理的财政和主权危机中。台湾的抵抗与日军登陆之间存在一段具体而暴力的地方过程，不能由条约签署取代。赔款、借款、租借地和铁路议题随即把中央财政同全球资本和港口网络接合；它们既扩大了技术、信息和商业的渠道，也加重了对税源、土地和地方合作的争夺。

1898年变法与政变显示改革并没有单一阵营或同一速度。诏令可以建立新的官制、教育和经济语言，不能证明省县已有经费、人员和执行程序。义和团运动、1900年北京陷落和1901年条约进一步说明宗教冲突、地方保护、外交武力与国家决策不能被一个标签解释。官府的“拳匪”、教会的“迫害”、列强的“保护”都是行动中的语言，而不是完整因果。

新政把学校、军队、警察、法律、商政和地方自治带入同一轮制度试验。1905年废科举重写了读书和入仕的资格，但私塾、家庭资源与地域不平等不会因此消失；新军和新式学堂也提供政治组织的空间，却不预定每位学生或士兵会走向革命。改革既增添国家能力，也暴露它需要省级财政、专业人员和社会信任才能延续。

\subsection{材料与争议}

上谕、章程、预算、报刊与外交照会应分别作为制度宣布、资金安排、公共争论和谈判立场来读。学校名册、警察档案、公司记录、地方诉讼和家庭材料只能局部检验实践。将改革描述为必然的现代化，或将其描述为全然失败，都会忽略参与扩大与门槛并存的事实。
